\documentclass[10pt,a4paper]{amsart}
\usepackage[margin=19mm]{geometry}
\usepackage{amsmath,amssymb,mathtools}
\usepackage[scr=rsfs,cal=euler]{mathalfa}
\usepackage{xfrac}
\usepackage[unicode,hidelinks,bookmarks=false]{hyperref}
\newcommand{\ie}{i.e.\ }
\newcommand{\cf}{cf.\ }
\newcommand{\resp}{respectively\ }
\newcommand{\Subsection}{\subsection}
\providecommand{\nobreakdash}{\nobreak\hbox{-}}
\setlength{\emergencystretch}{2em}
\multlinegap0pt
\usepackage{mathtools}
\usepackage{amssymb}
\usepackage[shortlabels]{enumitem}
\usepackage{float}
\usepackage{xcolor}
\usepackage{tikz}
\usepackage{algorithm}\usepackage{algpseudocode}
\usepackage{cancel}
\newtheorem{lemma}{Lemma}
\newtheorem{theorem}{Theorem}
\title{On Some Generalisations of Gauss Sequences}
\author{Sathyanarayan Narayan, N. Uday Kiran}
\date{Source: arXiv:2511.19503v5 (CC BY 4.0)}
\begin{document}
\maketitle

\noindent\emph{Source and licence.} On Some Generalisations of Gauss Sequences. Version arXiv:2511.19503v5 (CC BY 4.0), distributed by arXiv under the Creative Commons Attribution 4.0 International licence, http://creativecommons.org/licenses/by/4.0/, as recorded in the arXiv licence metadata for this exact version and shown on the abstract page https://arxiv.org/abs/2511.19503v5. Full licence notice: \texttt{licenses/arxiv-2511.19503-CC-BY-4.0.txt}.\par\smallskip

\noindent\textit{Editorial scope.} Counted unit: the article's theorem thm:1 (source0.tex lines 807--809). The article's complete original proof of it is delivered with it (source0.tex lines 810--849, one \texttt{proof} environment of 40 lines). No mathematics is written, repaired or re-derived: the statement and the proof are the article's own lines, in the article's own notation, and no retained line is substituted or re-flowed.  Retained uncounted as the source-local context the proof needs: lines 347--349; lines 394--396; lines 793--798.  Nothing was omitted from any retained span, so the package carries no omission record.  No retained line contains a citation, so the wrapper carries no bibliography and no bibliography entry.  No retained line contains a figure environment, an image-inclusion command or drawing code, so no image is delivered and the package declares no asset dependency.  Editorial formatting only, in addition: an AMS \texttt{amsart} preamble that restates the article's own declarations and macros used by the retained lines, namely the environments \texttt{lemma}, \texttt{theorem} on separate counters, together with the standard packages those lines need.  The retained environments print in this document as Lemma 1, Theorem 1 in order of appearance, whereas the article prints the same items with the numbering of its own full document.  The complete native source of the article as distributed in the arXiv e-print for this exact version is bundled unchanged as \texttt{sources/189723/source0.tex}.
\begin{equation}\label{eq:022}
a_n(q)=\sum_{d\mid n}\left[\frac{n}{d}\right]_{q^d}g_{\frac{n}{d}}(q^d)
\end{equation}

\begin{equation}\label{eq:011}
    a_n(\omega_n^j)=a_{\gcd(n,j)}(1)
\end{equation}

\begin{lemma}\label{lem:02}
    If a sequence $(a_n(q))\subset\mathbb{Z}[q]$ satisfies the $q$-Euler--Gauss congruence at a fixed $N\geq1$ and the $q$-Gauss congruence for all $1\leq n<N$, then 
    \[
        [N]_q\text{ divides } \sum_{d\mid N} \mu(d)a_{\frac{N}{d}}(q^d)\cdot\prod_{\substack{d\mid N, d>1\\\mu(d)=1}} a_{\frac{N}{d}}(q^d)
    \]
\end{lemma}

\begin{theorem}\label{thm:1} 
    Let $(a_n(q))$ be a $q$-Euler--Gauss sequence. If $a_{n}(1)\neq0$ for all $n$, then it is a $q$-Gauss sequence.
\end{theorem}

\begin{proof}
We prove the contrapositive. Suppose that $(a_n(q))$ is not a $q$-Gauss sequence, and let $n\geq1$ be the least positive integer at which the $q$-Gauss congruence fails. Note that $n$ cannot be a prime power (including $1$) because $(a_n(q))$ is a $q$-Euler--Gauss sequence. Now, from Lemma \ref{lem:02}, there exist divisors $D,d'>1$ of $n$ such that $\phi_D(q)$ divides $a_\frac{n}{d'}(q^{d'})$ where $\phi_D(q)$ denotes the $D$\textsuperscript{th} cyclotomic polynomial. We show that there exists $k\geq1$ such that $a_k(1)=0$, namely, 
\[a_\frac{n}{\operatorname{lcm}(d',D)}(1)=0.\] 
As $\{a_k(q)\}_{1\leq k<n}$ satisfies the $q$-Gauss congruence at the index $n/d'<n$, by the characterisation \eqref{eq:022} of $q$-Gauss sequences there exist $\{g_k(q)\}_{1\leq k< n}\subset\mathbb{Z}[q]$ such that, 
\[
    a_\frac{n}{d'}(q)=\sum_{d\mid\frac{n}{d'}}\left[\frac{n}{d'd}\right]_{q^{d}}g_\frac{n}{d'd}(q^{d})\text{ .}
\]
Since $\phi_D(q)\mid a_\frac{n}{d'}(q^{d'})$, the value of $a_\frac{n}{d'}(q)$ at the ${d'}$\textsuperscript{th} power of the $D$\textsuperscript{th} root of unity $\omega_D=\exp{\frac{2\pi i}{D}}$ is zero. On the other hand, evaluating the above expression at $q=\omega_D^{d'}$, and writing $e=d'd$, gives
\[
a_{\frac{n}{d'}}(\omega_D^{d'})
=
\sum_{\substack{e\mid n\\d'\mid e}}
\left[\frac{n}{e}\right]_{\omega_D^e}
g_{\frac{n}{e}}(\omega_D^e).
\]
Now, only those divisors $e\mid n$ divisible by both $d'$ and $D$, equivalently by $L:=\operatorname{lcm}(d',D)$, contribute to the sum because
\[\left[\frac{n}{e}\right]_{\omega_D^e}=
\begin{cases}
    \frac{n}{e}&D\mid e,\\
    0&D\nmid e
\end{cases}
\]
Consequently, writing $e=Le'$ we obtain
\[
a_{\frac{n}{d'}}(\omega_D^{d'})
=
\sum_{e'\mid\frac{n}{L}}
\frac{n}{Le'}g_{\frac{n}{Le'}}(1),
\]
since $\omega_D^e=1$ whenever $L\mid e$. By the characterisation \eqref{eq:022}, applied at the index $n/L$, the right-hand side is precisely $a_{n/L}(1)$. Therefore
\[
0=a_{\frac{n}{d'}}(\omega_D^{d'})=a_{\frac{n}{L}}(1),
\]
as claimed.


Alternatively, $\omega_D^{d'}$ can be rewritten as $\omega_\frac{n}{d'}^\frac{n}{D}$ and we may apply the characterisation of $q$-Gauss sequences \eqref{eq:011} at the index $n/d'$ so that 
\[0=a_\frac{n}{d'}(\omega_D^{d'})=a_\frac{n}{d'}(\omega_\frac{n}{d'}^\frac{n}{D})=a_{\gcd(\frac{n}{d'},\frac{n}{D})}(1).\]
Since $\gcd(\frac{n}{d'},\frac{n}{D})$ is precisely $\frac{n}{\operatorname{lcm}(d',D)}$, we have shown that $a_\frac{n}{L}(1)=0$.
\end{proof}

\end{document}
